% Study source only. Not routed into the Volume IV LaTeX tree.
% Purpose: stripped Set Theory proof queue from Volume I.

\section*{Gate 1 Set Theory Statements}

\subsection*{Scope}

This is a proof-order source for rebuilding the set-theoretic toolkit.  It
contains only axioms, definitions, and theorem-like statements.  It stops before
functions and relations.

\subsection*{Primitive Language and ZFC Axioms}

\paragraph{Primitive: Set and Membership.}
The objects of discourse are sets.  Membership is a primitive binary relation
written \(x\in A\).

\paragraph{Logical Equality / Identity.}
Set theory is formulated in first-order logic with equality.  Equality is the
logical identity relation:
\[
  A=A,
  \qquad
  A=B\to(\varphi(A)\Longleftrightarrow \varphi(B))
\]
for every formula \(\varphi\) in which substitution is legitimate.

\paragraph{Axiom: Extensionality.}
\[
  \forall A\,\forall B,\qquad
  A=B \Longleftrightarrow \forall x\,(x\in A\Longleftrightarrow x\in B).
\]

\paragraph{Axiom: Empty Set.}
\[
  \exists E\,\forall x,\ x\notin E.
\]

\paragraph{Axiom: Pairing.}
\[
  \forall a\,\forall b\,\exists P\,\forall x,\qquad
  x\in P \Longleftrightarrow (x=a\lor x=b).
\]

\paragraph{Axiom: Union.}
\[
  \forall A\,\exists U\,\forall x,\qquad
  x\in U \Longleftrightarrow \exists B\,(B\in A\land x\in B).
\]

\paragraph{Axiom: Power Set.}
\[
  \forall A\,\exists P\,\forall S,\qquad
  S\in P \Longleftrightarrow S\subseteq A.
\]

\paragraph{Axiom Schema: Separation.}
For every formula \(\varphi(x)\),
\[
  \forall A\,\exists B\,\forall x,\qquad
  x\in B \Longleftrightarrow (x\in A\land \varphi(x)).
\]

\paragraph{Axiom: Infinity.}
\[
  \exists A\,
  \bigl(\varnothing\in A
  \land
  \forall x\,(x\in A\to x\cup\{x\}\in A)\bigr).
\]

\paragraph{Axiom Schema: Replacement.}
For every formula \(\varphi(x,y)\),
\[
  \forall A\,
  \left(
  (\forall x\in A\,\exists!y\,\varphi(x,y))
  \to
  \exists B\,\forall y\,
  (y\in B\Longleftrightarrow \exists x\in A\,\varphi(x,y))
  \right).
\]

\paragraph{Axiom: Foundation.}
\[
  \forall A,\qquad
  A\neq\varnothing
  \to
  \exists x\in A\,(x\cap A=\varnothing).
\]

\paragraph{Axiom: Choice.}
\[
  \forall A,\qquad
  (\forall B\in A,\ B\neq\varnothing)
  \to
  \exists f\,\forall B\in A,\ f(B)\in B.
\]

\subsection*{Empty Set and Extensional Equality}

\paragraph{Theorem: Existence and Uniqueness of the Empty Set.}
\[
  \exists!E\,\forall x,\ x\notin E.
\]

\paragraph{Definition: Empty Set.}
\[
  \varnothing
  \quad\text{is the unique set satisfying}\quad
  \forall x,\ x\notin\varnothing.
\]

\paragraph{Definition: Subset.}
\[
  A\subseteq B
  \Longleftrightarrow
  \forall x\,(x\in A\to x\in B).
\]

\paragraph{Definition: Proper Subset.}
\[
  A\subsetneq B
  \Longleftrightarrow
  A\subseteq B\land A\neq B.
\]

\paragraph{Definition: Set Equality.}
\[
  A=B
  \Longleftrightarrow
  \forall x\,(x\in A\Longleftrightarrow x\in B).
\]

\paragraph{Theorem: Extensional Equality Criteria.}
\[
  A=B
  \Longleftrightarrow
  \forall x\,(x\in A\Longleftrightarrow x\in B),
\]
and
\[
  A=B
  \Longleftrightarrow
  (A\subseteq B\land B\subseteq A).
\]

\subsection*{Pairing and Union}

\paragraph{Theorem: Existence and Uniqueness of Pairing's Output.}
\[
  \forall a\,\forall b\,\exists!P\,\forall x,\qquad
  x\in P\Longleftrightarrow (x=a\lor x=b).
\]

\paragraph{Definition: Pair and Singleton.}
\[
  \{a,b\}
  \quad\text{is the unique set whose members are exactly }a\text{ and }b,
\]
and
\[
  \{a\}:=\{a,a\}.
\]

\paragraph{Theorem: Existence and Uniqueness of Union's Output.}
\[
  \forall A\,\exists!U\,\forall x,\qquad
  x\in U\Longleftrightarrow \exists B\,(B\in A\land x\in B).
\]

\paragraph{Definition: Union Over a Set.}
\[
  \bigcup A
  \quad\text{is the unique set satisfying}\quad
  x\in\bigcup A
  \Longleftrightarrow
  \exists B\,(B\in A\land x\in B).
\]

\paragraph{Corollary: Existence and Uniqueness of Binary Union.}
\[
  \forall A\,\forall B\,\exists!C\,\forall x,\qquad
  x\in C\Longleftrightarrow (x\in A\lor x\in B).
\]

\paragraph{Definition: Binary Union.}
\[
  A\cup B:=\bigcup\{A,B\},
  \qquad
  x\in A\cup B
  \Longleftrightarrow
  (x\in A\lor x\in B).
\]

\paragraph{Theorem: Union is Inclusion-Monotone.}
If \(A\subseteq B\), then
\[
  A\cup C\subseteq B\cup C
  \quad\text{and}\quad
  C\cup A\subseteq C\cup B.
\]

\subsection*{Separation, Intersection, Difference, and Complement}

\paragraph{Theorem: Existence and Uniqueness of Separation's Output.}
For every formula \(\varphi(x)\),
\[
  \forall A\,\exists!B\,\forall x,\qquad
  x\in B\Longleftrightarrow (x\in A\land\varphi(x)).
\]

\paragraph{Corollary: Existence and Uniqueness of Intersection.}
\[
  \forall A\,\forall B\,\exists!C\,\forall x,\qquad
  x\in C\Longleftrightarrow (x\in A\land x\in B).
\]

\paragraph{Definition: Intersection.}
\[
  A\cap B:=\{x\in A\mid x\in B\},
  \qquad
  x\in A\cap B
  \Longleftrightarrow
  (x\in A\land x\in B).
\]

\paragraph{Theorem: Intersection is Inclusion-Monotone.}
If \(A\subseteq B\), then
\[
  A\cap C\subseteq B\cap C
  \quad\text{and}\quad
  C\cap A\subseteq C\cap B.
\]

\paragraph{Corollary: Existence and Uniqueness of Set Difference.}
\[
  \forall A\,\forall B\,\exists!C\,\forall x,\qquad
  x\in C\Longleftrightarrow (x\in A\land x\notin B).
\]

\paragraph{Definition: Set Difference.}
\[
  A\setminus B:=\{x\in A\mid x\notin B\},
  \qquad
  x\in A\setminus B
  \Longleftrightarrow
  (x\in A\land x\notin B).
\]

\paragraph{Theorem: Set Difference is Inclusion-Monotone in the Left Argument.}
If \(A\subseteq B\), then
\[
  A\setminus C\subseteq B\setminus C.
\]

\paragraph{Theorem: Set Difference is Inclusion-Antitone in the Right Argument.}
If \(C\subseteq D\), then
\[
  A\setminus D\subseteq A\setminus C.
\]

\paragraph{Corollary: Existence and Uniqueness of Relative Complement.}
For a fixed ambient set \(U\) and every set \(A\),
\[
  \exists!C\,\forall x,\qquad
  x\in C\Longleftrightarrow (x\in U\land x\notin A).
\]

\paragraph{Definition: Complement Relative to \(U\).}
\[
  A^c:=U\setminus A,
  \qquad
  x\in A^c
  \Longleftrightarrow
  (x\in U\land x\notin A).
\]

\paragraph{Theorem: Complement is Inclusion-Antitone.}
If \(A\subseteq B\subseteq U\), then
\[
  B^c\subseteq A^c.
\]

\paragraph{Corollary: Existence and Uniqueness of Symmetric Difference.}
\[
  \forall A\,\forall B\,\exists!C\,\forall x,\qquad
  x\in C
  \Longleftrightarrow
  ((x\in A\land x\notin B)\lor(x\in B\land x\notin A)).
\]

\paragraph{Definition: Symmetric Difference.}
\[
  A\triangle B:=(A\setminus B)\cup(B\setminus A),
\]
so
\[
  x\in A\triangle B
  \Longleftrightarrow
  ((x\in A\land x\notin B)\lor(x\in B\land x\notin A)).
\]

\subsection*{Power Sets}

\paragraph{Theorem: Existence and Uniqueness of Power Set's Output.}
\[
  \forall A\,\exists!P\,\forall S,\qquad
  S\in P\Longleftrightarrow S\subseteq A.
\]

\paragraph{Definition: Power Set.}
\[
  \mathcal P(A)
  \quad\text{is the unique set satisfying}\quad
  S\in\mathcal P(A)\Longleftrightarrow S\subseteq A.
\]

\paragraph{Theorem: Power Set is Inclusion-Monotone.}
If \(A\subseteq B\), then
\[
  \mathcal P(A)\subseteq\mathcal P(B).
\]

\subsection*{Boolean Set Laws}

\paragraph{Theorem: De Morgan's Laws.}
For \(A,B\subseteq U\),
\[
  (A\cup B)^c=A^c\cap B^c
  \quad\text{and}\quad
  (A\cap B)^c=A^c\cup B^c.
\]

\paragraph{Theorem: Commutativity.}
\[
  A\cup B=B\cup A,
  \qquad
  A\cap B=B\cap A.
\]

\paragraph{Theorem: Associativity.}
\[
  (A\cup B)\cup C=A\cup(B\cup C),
  \qquad
  (A\cap B)\cap C=A\cap(B\cap C).
\]

\paragraph{Theorem: Distributivity.}
\[
  A\cap(B\cup C)=(A\cap B)\cup(A\cap C),
\]
and
\[
  A\cup(B\cap C)=(A\cup B)\cap(A\cup C).
\]

\paragraph{Theorem: Identity and Absorption Laws.}
\[
  A\cup\varnothing=A,
  \qquad
  A\cap U=A,
\]
and
\[
  A\cup(A\cap B)=A,
  \qquad
  A\cap(A\cup B)=A.
\]

\paragraph{Theorem: Idempotency.}
\[
  A\cup A=A,
  \qquad
  A\cap A=A.
\]

\paragraph{Theorem: Complement Laws.}
\[
  A\cup A^c=U,
  \qquad
  A\cap A^c=\varnothing,
\]
and
\[
  \varnothing^c=U,
  \qquad
  U^c=\varnothing.
\]

\paragraph{Theorem: Involution.}
\[
  (A^c)^c=A.
\]

\paragraph{Theorem: Difference Laws.}
\[
  A\setminus B=A\cap B^c,
\]
\[
  A\setminus\varnothing=A,
  \qquad
  A\setminus A=\varnothing,
  \qquad
  \varnothing\setminus A=\varnothing,
\]
and
\[
  A\setminus(B\cup C)=(A\setminus B)\cap(A\setminus C),
  \qquad
  A\setminus(B\cap C)=(A\setminus B)\cup(A\setminus C).
\]

\paragraph{Theorem: Symmetric Difference Laws.}
\[
  A\triangle B=B\triangle A,
  \qquad
  A\triangle\varnothing=A,
  \qquad
  A\triangle A=\varnothing.
\]
Also,
\[
  A\triangle B=\varnothing\Longleftrightarrow A=B.
\]

\paragraph{Theorem: Subset Absorption Criteria.}
\[
  A\subseteq B
  \Longleftrightarrow
  A\cup B=B,
\]
and
\[
  A\subseteq B
  \Longleftrightarrow
  A\cap B=A.
\]

\subsection*{Indexed Families of Sets}

\paragraph{Definition: Indexed Family of Sets.}
An indexed family of sets is a collection \(\{A_i\}_{i\in I}\) assigning a set
\(A_i\) to each index \(i\in I\).

\paragraph{Definition: Indexed Union.}
\[
  \bigcup_{i\in I}A_i
  :=
  \{x\mid \exists i\in I,\ x\in A_i\},
\]
so
\[
  x\in\bigcup_{i\in I}A_i
  \Longleftrightarrow
  \exists i\in I,\ x\in A_i.
\]

\paragraph{Definition: Indexed Intersection.}
\[
  \bigcap_{i\in I}A_i
  :=
  \{x\mid \forall i\in I,\ x\in A_i\},
\]
so
\[
  x\in\bigcap_{i\in I}A_i
  \Longleftrightarrow
  \forall i\in I,\ x\in A_i.
\]

\paragraph{Theorem: Indexed De Morgan Laws.}
\[
  U\setminus\bigcup_{i\in I}A_i
  =
  \bigcap_{i\in I}(U\setminus A_i).
\]
If \(I\neq\varnothing\), then
\[
  U\setminus\bigcap_{i\in I}A_i
  =
  \bigcup_{i\in I}(U\setminus A_i).
\]

\paragraph{Theorem: Indexed Distributivity.}
\[
  A\cap\bigcup_{i\in I}B_i
  =
  \bigcup_{i\in I}(A\cap B_i),
\]
and
\[
  A\cup\bigcap_{i\in I}B_i
  =
  \bigcap_{i\in I}(A\cup B_i),
\]
with the usual nonempty-index convention for indexed intersections.

\subsection*{Stop Point}

Stop here before functions, relations, images, preimages, order relations,
set-family closure systems, set algebras, sigma-algebras, topology, or Borel
generation.
